%%% main.tex --- Elenchon: AI-Assisted Requirements-Based Testing via Cause-Effect Graphs
%%%
%%% Portable skeleton: compiles with a stock TeX Live (base `article` + common
%%% packages). Build:  latexmk -pdf main.tex   (or `make`).

\documentclass[11pt,a4paper]{article}

\usepackage[utf8]{inputenc}
\usepackage[T1]{fontenc}
\usepackage{lmodern}
\usepackage[margin=1in]{geometry}
\usepackage{microtype}
\usepackage{booktabs}
\usepackage{listings}
\usepackage{xcolor}
\usepackage[numbers,sort&compress]{natbib}
\usepackage[hidelinks]{hyperref}

\lstset{basicstyle=\ttfamily\small, columns=fullflexible, keepspaces=true,
  showstringspaces=false, frame=single, framesep=4pt, breaklines=true}

\title{\textbf{Elenchon}: AI-Assisted Requirements-Based\\ Testing via Cause-Effect Graphs}
\author{Bob\\ \small\texttt{eternal.recursion@proton.me}}
\date{\today}

\begin{document}
\maketitle

\begin{abstract}
Most software defects originate in requirements, not code, yet requirements are
written in ambiguous natural language and rarely tested as first-class artifacts.
Cause-effect graphing---formalized industrially as Richard Bender's
Requirements-Based Testing---addresses this by reducing a requirement to a boolean
graph of \emph{causes} and \emph{effects} from which a minimal, high-coverage set
of functional test cases can be derived. Its adoption has been limited by the
manual labor of the reduction. We present \textsc{Elenchon}, a statically typed
(\textsc{Coalton}) cause-effect-graph engine: given a \emph{structured} graph of
causes and effects it reasons to the minimum decision table demonstrating each
cause independently drives its effect (essentially MC/DC coverage), and it surfaces
requirement defects---ambiguity, incompleteness, contradiction---as a co-product.
Closing the manual-reduction gap is the job of a \emph{consumer above} Elenchon: an
agentic, large-language-model front-end (the \textsc{ChatRBT} example built on
\textsc{praxeon}) that \emph{disambiguates and formalizes} prose requirements into
the structured causes and effects Elenchon consumes. That front-end depends on
Elenchon, not the reverse; ambiguity is handled there not by guessing but through
the Common Lisp condition system, which escalates un-formalizable requirements back
to the human. Elenchon is the \emph{fourth} framework in a six-framework ecosystem
whose thesis is to establish Coalton as a first-class ``modern Lisp''; it depends
only leftward---on the functional-collections library (\textsc{aion}) for the graph
and decision-table data---and never on the agentic layer that drives it.
\end{abstract}

\section{Introduction}
\label{sec:intro}
% Defects are born in requirements; requirements are ambiguous and untested; CEG/RBT
% as the corrective; the manual-reduction bottleneck; LLMs as the missing front-end.

\section{Background: Cause-Effect Graphing and RBT}
\label{sec:background}
% Myers' technique; Bender's RBT industrialization; causes/effects; the E/I/O/R/M
% constraints; decision tables; the minimal-set / MC/DC coverage criterion.

\section{The Formalization Gap}
\label{sec:gap}
% Why CEG/RBT is underused: turning prose into causes/effects is skilled manual
% labor. The disambiguation problem stated precisely.

\section{Disambiguation as an LLM Task Under Rigor (the Consumer Layer)}
\label{sec:extract}
% This layer lives ABOVE Elenchon, in a consumer (praxeon's ChatRBT) that depends on
% Elenchon and feeds it structured causes/effects -- Elenchon itself does no NL.
% NL -> yes/no causes & effects; constraining LLM output to propositions; ambiguity
% as a recoverable failure via the condition system (retry/substitute/abandon->ask);
% propose-and-confirm vs autonomous. Realized as a praxeon End/Means/Action loop.

\section{A Typed Cause-Effect-Graph Core}
\label{sec:ceg}
% The CEG as an algebraic data type in Coalton; boolean wiring + constraints;
% expression-tree vs node-graph representation; why static types fit this core.

\section{Solving for Minimal, Sufficient Tests}
\label{sec:solve}
% CEG -> decision table; the minimum columns showing each cause independently drives
% its effect; native Bender/Myers table-walk vs an FFI'd SAT/SMT backend behind a
% neutral protocol; traceability as a first-class output.

\section{Requirement Defects as a Co-Product}
\label{sec:defects}
% Dangling causes/effects, contradictory constraints, non-propositional requirements;
% the ambiguity review as arguably the higher-value deliverable.

\section{The Ecosystem}
\label{sec:ecosystem}
% The six core frameworks in dependency order (low->high): aion (functional stdlib),
% cons (project & dev tooling), mnemosyne (bitemporal data layer), elenchon (this,
% the CEG/RBT engine -- the fourth), hyperion (HTMX web framework), praxeon (agentic
% AI), plus hermes (a satellite leaf-lib for external integrations -- email/SMS,
% payments later -- outside the linear DAG; it depends only on aion, and nothing in
% the core DAG depends on it).
% A framework never depends on one to its right. The shared house style; Coalton
% as a first-class modern Lisp. Elenchon depends only leftward (aion; later mnemosyne)
% and never on praxeon/hyperion; the agentic NL->CEG layer (praxeon's ChatRBT) is a
% consumer ABOVE Elenchon that depends on it, not the reverse.

\section{Evaluation Plan}
\label{sec:eval}
% Test-count and coverage vs hand-written suites and, where available, Bender RBT
% output; extraction fidelity (human-confirmed causes/effects); a worked canonical
% example (requirement -> graph -> table -> tests).

\section{Related Work}
\label{sec:related}
% Bender RBT/SoftTest; Myers cause-effect graphing; MC/DC (DO-178C); combinatorial
% (pairwise) and classification-tree test design; LLM-based test generation.

\section{Conclusion}
\label{sec:conclusion}

\bibliographystyle{plainnat}
\bibliography{references}

\end{document}
